% ===========================================================================
\section{Introduction}
\label{sec:introduction}
% ===========================================================================

Waste is the category of a national greenhouse-gas inventory where the arithmetic is easiest
and the activity data hardest. The emission factors are few and well characterised; what is
difficult is knowing how much waste a country generates, what fraction of it reaches a
disposal site, what happens to the fraction that does not, and how the sewage of a large
population is actually handled. None of those quantities is metered anywhere in Pakistan.

That difficulty has a specific consequence for how the sector is usually compiled. Because the
activity data must be modelled rather than read, each category tends to be estimated on its
own terms: landfill methane from a tonnage assumption, open burning from a per-capita burning
rate in the literature, wastewater from a population and a default organic load. The three
estimates are then reported together. Nothing in that procedure requires the waste sent to
landfill and the waste openly burned to sum to a quantity of waste the country plausibly
produces, and in general they do not.

\subsection{What this paper does}

We report Pakistan's waste sector for FY2023--24 from Pakistani activity data, with the three
categories constructed so that the constraint the conventional procedure lacks is enforced by
construction. A single national waste-generation figure is partitioned: what is collected
goes to the first-order-decay model for solid-waste disposal, what is not collected is what
can burn, and the two sum to the generated quantity exactly. The partition is asserted in code
with equality rather than checked afterwards with a tolerance, because a tolerance permits
precisely the drift the constraint exists to remove.

The result is a sector total of \SI{37.45}{\mega\tonne} \COe{}, of which domestic wastewater
is two-thirds. It is the second most uncertain sector in the national inventory that contains
it, and this paper is as much about where that uncertainty lives and what would reduce it as
about the central estimate.

\subsection{Why Pakistan, and why this sector now}

Pakistan is the fifth most populous country in the world, urbanising rapidly, and its most
recent full national inventory submission reports for FY2017--18 \citep{pakbur1}. In the
interim the international gridded products---EDGAR \citep{edgar2025booklet,crippa2020} and Climate
TRACE \citep{climatetrace}---supply the numbers that inform policy. They are global models
applying globally consistent assumptions to nationally reported activity, which is their
strength for cross-country comparison and their limitation for any one country: where a
national compiler can consult a sanitation survey or a district population panel, a global
product must choose a default.

The waste sector is where that difference bites hardest, for two reasons. Its activity data
are the least standardised internationally, so the defaults do the most work. And it is
dominated by methane, whose short atmospheric lifetime makes it the gas where near-term
mitigation is most effective and therefore where the estimate most directly changes what a
government should do.

\subsection{Contributions}

\begin{enumerate}[leftmargin=*,itemsep=2pt,topsep=3pt]
\item \textbf{A current national waste-sector inventory for Pakistan} covering solid-waste
  disposal, open burning and domestic wastewater, from a district population panel and
  national sanitation statistics rather than from a global model's Pakistani slice.
\item \textbf{One waste stream, mass-balanced across its fates}: collected and uncollected
  shares sum to generation exactly, two fates carry emissions pathways (4A, 4C2) and the
  unburned remainder ships as an explicit unresolved row rather than a silent omission. The
  closed balance removes a class of inconsistency that independent estimation cannot
  detect (Section~\ref{sec:methods}).
\item \textbf{Both readings of the sewer methane correction factor}, with the evidence for
  each stated and neither resolved by fiat, together with the measurement that would settle
  it (Section~\ref{sec:discussion}).
\item \textbf{A spin-up sensitivity for the decay model}, quantifying how much of today's
  landfill methane comes from waste deposited before the activity panel begins---a question
  most first-order-decay implementations answer implicitly and do not report.
\item \textbf{An air-pollutant layer on the same activity basis} as the greenhouse gases, so
  that open burning's particulate can be compared with other Pakistani sources without the
  usual mismatched-activity problem.
\end{enumerate}

Method that this sector shares with the other nine of the national inventory---the module
contract, the uncertainty treatment, the notation keys, the accounting closures---is
documented in the national manuscript \citep{ahmad2026national} and not repeated here.
